\chapter{Numerical setup and implementation}
\label{chap:numerical-setup}

\section{Code path}
The implementation is organized around the program \texttt{src/test\_laplace\_RT\_scan.c} and the accompanying analysis scripts in \texttt{laplace\_static/scripts}. The production code path is:
\begin{enumerate}
  \item read gauge fields and precomputed Laplacian eigenvectors from the existing QCD-library/HDF5 setup;
  \item build static perambulators \(\tau_{ij}(\bm x;t_0,t_1)\) for the selected time extent \(\tau\);
  \item contract pairs of static perambulators to form \(L(R,\tau)\);
  \item average over spatial positions, equivalent orientations, configurations, and selected temporal source positions;
  \item analyze effective energies, plateau windows, potential fits, and reliability diagnostics.
\end{enumerate}

\section{Current production parameters}
The current benchmark setup is summarized in \cref{tab:production-setup}. These values reflect the validated production stage at which the thesis write-up is being started.

\begin{table}[htbp]
\centering
\caption{Current production and diagnostic parameters for the Laplace-static benchmark.}
\label{tab:production-setup}
\begin{tabular}{ll}
\toprule
Quantity & Current value \\
\midrule
Gauge configurations & \(N_{\mathrm{cfg}}=100\) \\
Eigenvectors & \(N_v=10\) \\
Eigenmode profile & constant, \(\rho_i=1\) \\
Temporal-source average & \(t_0/a=0,8,16,24,32,40\) \\
Validated scan & \(R/a=1,\ldots,6\), \(\tau/a=1,\ldots,8\) \\
Diagnostic scan & \(R/a=1,\ldots,12\), \(\tau/a=1,\ldots,10\) \\
Stable benchmark window & \(\tau/a=3,4,5\) \\
\bottomrule
\end{tabular}
\end{table}

\section{Statistical analysis}
The analysis uses correlated fits where possible, with jackknife and Gamma-method style cross-checks. The Gamma method estimates autocorrelation effects through explicit autocorrelation functions and an automatically chosen summation window \cite{Wolff2004}. For derived observables and fit quantities, the thesis should distinguish clearly between the central estimator, the covariance estimate, the fit-window choice, and the treatment of autocorrelations.

\todo{Once the final analysis scripts are frozen, document exact commands, file names, software versions, and random/fit settings in this chapter and in \cref{app:reproducibility}.}
